%% kapa: exam=1 round=35 year=2024 period=1 form=A subject=real_estate subject_ko=부동산학원론 seq=27 number=107
%% source: Q-Net 감정평가사 기출문제 (cst00309) artlSeq=5213229 file=raw/1008/5213229/1교시 문제지 원본.pdf
%% conversion: kordoc 4.24.51 -> tools/gen1.py (auto); review=proofread-sonnet; flags=none
\begin{question}{107}
아파트시장의 균형가격과 균형거래량에 관한 설명으로 옳지 않은 것은? (단, 완전탄력적과 완전비탄력적 조건이 없는 경우는 수요와 공급의 법칙에 따르며, 다른 조건은 동일함)
\begin{choices}
  \item 수요의 증가폭이 공급의 증가폭보다 클 경우, 균형가격은 하락하고 균형거래량은 증가한다.
  \item 균형상태인 아파트시장에서 건축원자재의 가격이 상승하면 균형가격은 상승하고 균형거래량은 감소한다.
  \item 공급이 가격에 대해 완전탄력적인 경우, 수요가 증가하면 균형가격은 변하지 않고 균형거래량만 증가한다.
  \item 공급이 가격에 대해 완전비탄력적인 경우, 수요가 증가하면 균형가격은 상승하고 균형거래량은 변하지 않는다.
  \item 공급의 감소폭이 수요의 감소폭보다 클 경우, 균형가격은 상승하고 균형거래량은 감소한다.
\end{choices}
\end{question}
\answer{1}
